\paragraph{C.3at (Refuted, with specified counterexamples): ordinary restriction fails at $N=r+2$ in every positive characteristic.}
For every prime $p$ and integer $r\ge2$, the construction below specifies
an algebraically closed field $k$ of characteristic $p$, a regular
excellent strictly henselian local $k$-ring $R$ of dimension $N=r+2$,
and an element $f\in R$ for which the two punctures
\[
 X=\operatorname{Spec}R\setminus\{\mathfrak m_R\},
 \qquad Y=X\cap V(f)
\]
are regular and
\[
 \mathrm{CH}^r(X)=0,
 \qquad
 \mathrm{CH}^r(Y)_{\mathrm{alg}}\otimes\mathbb Q\ne0.
\]
Thus ordinary $J^r$ restriction is not universally surjective at
$N=r+2$, even rationally, in any positive characteristic. The field is
explicitly a transcendental extension of $\overline{\mathbb F}_p$:
\[
             k=\overline{\overline{\mathbb F}_p(T_0)}.
\]
There is no assertion of rational nonvanishing over
$\overline{\mathbb F}_p$ itself, where ordinary algebraically trivial
cycles are torsion by C.3o.

\begin{proof}
\emph{Step 1: a specified smooth Calabi--Yau hypersurface with nonzero
Hasse coefficient.}
Put $m=r+2$ and use $m$ variables $x_0,\ldots,x_{m-1}$. Over
$\mathbb F_p$ define
\[
 F_0=
 \begin{cases}
  \sum_{i=0}^{m-1}x_i^m,&p\nmid m,\\[2pt]
  \displaystyle\sum_{i=0}^{m-2}x_i^{m-1}x_{i+1}+x_{m-1}^m,
       &p\mid m,
 \end{cases}
 \qquad
 F_\lambda=F_0+\lambda\prod_{i=0}^{m-1}x_i.
\]
The hypersurface of $F_0$ is smooth. For $p\nmid m$ this follows
from the Fermat gradient. For $p\mid m$, the last partial derivative
is $x_{m-2}^{m-1}$; successive preceding derivatives force
$x_{m-2}=x_{m-3}=\cdots=x_0=0$ at a simultaneous zero of all partials.
The equation $F_0=0$ then forces $x_{m-1}=0$. Here $m\ge4$, so all
the exponents used in that induction are positive.

Let
\[
 b(\lambda)=
 [x_0^{p-1}\cdots x_{m-1}^{p-1}]F_\lambda^{p-1}
       \in\mathbb F_p[\lambda].
\]
Its leading term is $\lambda^{p-1}$, so it is nonzero. The bad
parameters for smoothness form a proper closed subset of the affine
parameter line: the relative singular locus is closed in
$\mathbb P^{m-1}\times\mathbb A^1$, its projection is closed by
properness, and it omits $\lambda=0$. Thus only finitely many
parameters are forbidden by smoothness or $b(\lambda)=0$.

Here is a terminating finite-field rule specifying a permissible
parameter, rather than an unspecified generic choice. Order the monic
irreducible polynomials $g(z)\in\mathbb F_p[z]$ first by degree and
then lexicographically by their coefficient vectors, representing field
elements by $0,\ldots,p-1$. For each $g$, put
$\mathbb F_q=\mathbb F_p[z]/(g)$ and $\lambda=z\bmod g$.
Choose the first $g$ for which $b(\lambda)\ne0$ and the ideal of
$F_\lambda$ and all its partial derivatives is the unit ideal in every
projective affine chart $x_i=1$. These are finite polynomial calculations.
The rule terminates because the nonempty open just proved has a point
over a finite extension of $\mathbb F_p$. Define
\[
 T_0=V(F_\lambda)\subset\mathbb P^{r+1}_{\mathbb F_q},
 \qquad k_0=\overline{\mathbb F}_p,
 \qquad T_*=T_0\times_{\mathbb F_q}k_0.
\]
This is a smooth geometrically integral hypersurface of dimension $r$.
Indeed a smooth hypersurface of positive dimension here is geometrically
connected by the standard projective-space cohomology sequence; its
smooth irreducible components would otherwise be disjoint.
Adjunction and the same sequence give
\[
 K_{T_0}\simeq\mathcal O_{T_0},\qquad
 H^i(T_0,\mathcal O_{T_0})=0\ (0<i<r),\qquad
 H^r(T_0,\mathcal O_{T_0})\simeq\mathbb F_q.
\]

\emph{Step 2: the slope-zero class, including $p=2$ and $p\mid m$.}
The connecting isomorphism
\[
 H^r(T_0,\mathcal O_{T_0})
  \simeq H^{r+1}(\mathbb P^{r+1},\mathcal O(-m))
\]
identifies its generator with the Cech class
$(x_0\cdots x_{m-1})^{-1}$. Under absolute Frobenius, the connecting
diagram multiplies the $p$th power of this representative by
$F_\lambda^{p-1}$. In the top Cech quotient the only surviving monomial
of total degree $-m$ has every exponent $-1$. Frobenius on the
one-dimensional cohomology group is therefore
\[
                         c\longmapsto b(\lambda)c^p.
\]
This is bijective. This Cech computation is also the projective-space
case of the Hasse--Witt formula of
\citet[Section~2, Corollaries 2.3 and 2.5, author pp.~8--9]{D_HuangLianYauYuHasse}.
No condition $p\nmid m$ is imposed in this computation.

For clarity, the passage from this coherent cohomology computation to
nonzero rational crystalline cohomology is proved here. Write
$W=W(\mathbb F_q)$, let $\sigma$ denote its Frobenius automorphism, and set
$M_n=H^r(T_0,W_n\mathcal O_{T_0})$. The Witt exact sequence and
$H^{r-1}(T_0,\mathcal O)=0$ give
\[
 0\longrightarrow M_{n-1}\xrightarrow{V_n}M_n
       \longrightarrow H^r(T_0,\mathcal O)\longrightarrow0.
\]
Here $H^{r+1}(T_0,W_{n-1}\mathcal O)=0$ follows from the finite Witt
filtration and coherent cohomological dimension. Verschiebung $V_n$ is
$\sigma^{-1}$-semilinear as a $W$-module map; this does not change
finite length. Induction gives
$\operatorname{length}_W(M_n)=n$ and
$H^{r-1}(T_0,W_n\mathcal O)=0$. The coordinatewise $p$th-power
operator $F_n:M_n\to M_n$ is $\sigma$-semilinear; the identity
$F_nV_n=V_nF_{n-1}$ and induction on this short exact sequence show
that $F_n$ is bijective. Truncation
$R_n:M_n\to M_{n-1}$ is surjective: the kernel sheaf of
$W_n\mathcal O\to W_{n-1}\mathcal O$ is a Frobenius twist of
$\mathcal O$, whose cohomology in degree $r+1$ vanishes.
The correctly typed identity
$p\operatorname{id}_{M_n}=V_nF_{n-1}R_n$ now implies
\[
 pM_n=V_n M_{n-1},\qquad M_n/pM_n\simeq\mathbb F_q.
\]
Nakayama's lemma makes $M_n$ cyclic over $W_n(\mathbb F_q)$; its
length identifies it with that ring. Choose compatible generators using
surjective truncation and its compatibility with projection to
$H^r(\mathcal O)$. We obtain
\[
 H^r(T_0,W\mathcal O_{T_0})
       =\varprojlim_n M_n\simeq W,
\]
with bijective Frobenius. The inverse-limit equality is the unnumbered
proper-cohomology comparison at the beginning of
\citet[Section~2, author pp.~6--7]{D_BerthelotBlochEsnault2007};
the Witt exact sequence is also the special case $I=\mathcal O$ of
their Lemma~2.5, equation~(2.7), author pp.~9--10.
Their Theorem~1.1, or the smooth proper comparison (1.1) recalled
in their introduction, identifies its rationalization with
\[
 E=H^r_{\mathrm{cris}}(T_0/W)[1/p]_{<1}.
\]
Consequently $E$ is nonzero; in fact it has rank one and slope zero.
This remains true after every finite extension of the coefficient
field. The argument proves exactly the slope needed below and does
not assume that all cohomological degrees of $T_0$ are ordinary.

\emph{Step 3: a generic-point cycle cannot be supported on a divisor.}
Set $K=k_0(T_*)$, $k=\overline K$, and $T=T_*\times_{k_0}k$.
Let $P\in T(k)$ be the geometric generic point specified by the
inclusion $K\subset k$. Put
\[
 H=c_1(\mathcal O_T(1)),\qquad
 D=T_*\cap V(x_{m-1}),
\]
and choose a constant point $a\in D(k_0)$. Such a point may be
specified by enumerating projective points over finite extensions until
$F_\lambda=x_{m-1}=0$; this terminates because $D$ is a nonempty
projective hypersurface of dimension $r-1$. Define
$\alpha=[P]-[a]\in\mathrm{CH}_0(T)$.

Suppose a nonzero multiple of $\alpha$ lay in
$H\mathrm{CH}^{r-1}(T)_{\mathbb Q}$. Every hyperplane product has a
representative supported on $D_k$. After clearing denominators the
supposed rational equivalence is finite data, so it descends from
$k=\overline K$ to a finite extension $K'/K$. Proper pushforward
along $T_{K'}\to T_K$ gives
\[
                 b[\eta]=z_D\quad\text{in }\mathrm{CH}_0(T_K),
                 \qquad b\ne0,
\]
where $z_D$ is supported on $D_K$. The constant $a$ is absorbed
into that support. Inseparable degrees cause no difficulty: the norm
multiplies the generic point by the positive integer $[K':K]$, and
we are working rationally. Spreading from the first generic factor
and applying Chow localization yields
\[
 b[\Delta_{T_*}]=\Gamma_D+\Gamma_Z
       \quad\text{in }\mathrm{CH}^r(T_*\times T_*)_{\mathbb Q},
\]
where $\Gamma_D$ is supported on $T_*\times D$ and $\Gamma_Z$
on $Z\times T_*$ for a proper divisor $Z\subset T_*$. The norm,
spreading and diagonal-localization construction is exactly the one
recalled in \citet[Section~1, printed p.~187]{D_Esnault2003}.
All cycles, rational functions and closed supports in this equality
descend to one finite extension of $\mathbb F_q$.

We explain both correspondence actions on the nonzero slope-$<1$
space $E$ after that extension. The action of $\Gamma_D$ lands in
the image of supported rigid cohomology
$H^r_D(T_*/\operatorname{Frac}W)$, since its restriction to
$T_*\setminus D$ is zero. By
\citet[Lemma~2.1 and its proof, printed pp.~189--190]{D_Esnault2003},
every slope of this supported group is at least one. Thus this
Frobenius-equivariant action is zero on $E$.

For $\Gamma_Z$ one must also justify the source support; differential
forms alone would not suffice in positive characteristic. For each
irreducible component of this correspondence, let $Z_1$ be its
source image and choose a smooth projective alteration
$Z'_1\to Z_1$, with $\dim Z'_1\le r-1$. Such an alteration exists
over a finite extension by
\citet[Theorem~4.1 and Remark~4.2, printed p.~66]{D_DeJong1996Alterations}.
The original component lifts to a rational cycle on $Z'_1\times T_*$:
over its generic point choose a closed point of the nonempty proper
fibre, take its closure, and divide by its finite generic degree.
Its action consequently factors through
$H^r_{\mathrm{cris}}(Z'_1)[1/p]$. By
\citet[Theorem~1.1 and Section~2]{D_BerthelotBlochEsnault2007},
the slope-$<1$ part of the latter group is
$H^r(Z'_1,W\mathcal O)[1/p]=0$, since $\dim Z'_1<r$.
The last vanishing follows either from Zariski cohomological dimension
and the stated inverse-limit comparison, or by the finite-length Witt
filtration and the same comparison. The pullback from $T_*$ preserves
Frobenius slopes, so $\Gamma_Z$ also kills $E$.

These cycle-class operations factor through rational Chow groups and
are compatible with pullback, proper pushforward and Frobenius, as
spelled out in \citet[Section~2, printed p.~189]{D_Esnault2003}.
Thus the diagonal equality would make $b\operatorname{id}_E=0$.
The nonzero integer $b$ is invertible in the characteristic-zero
coefficient field, even when $p\mid b$, a contradiction. Hence
\[
       \alpha\ne0\quad\text{in}\quad
       \mathrm{CH}^r(T)_{\mathbb Q}/
                     H\mathrm{CH}^{r-1}(T)_{\mathbb Q}.
\]

\emph{Step 4: actual ordinary equivalence and the local construction.}
The points $P,a$ lie on a smooth connected projective curve inside
$T$: take successive sufficiently ample hypersurfaces through the two
points, with independent prescribed first jets, and apply Bertini over
the infinite algebraically closed field $k$
\citep[Lemmas 33.47.1 and 33.47.3]{D_StacksBertini0FD4}. This uses a very ample
embedding and its jet-separated systems, so is valid in every
characteristic. Its diagonal correspondence followed by the inclusion
in $T$ is an actual family on a fixed product with a smooth proper
parameter curve. It proves $\alpha\in\mathrm{CH}^r(T)_{\mathrm{alg}}$.

The homogeneous hypersurface ring is normal and generated in degree
one: it is Cohen--Macaulay, and its puncture is smooth because $T$
is smooth. Projective-space cohomology identifies it with the full
section ring of $\mathcal O_T(1)$. Define
\[
 R=k[x_0,\ldots,x_{r+1}]_{(x_0,\ldots,x_{r+1})}^{\mathrm{sh}},
 \qquad f=F_\lambda\in R.
\]
Strict henselization is taken at the geometric point with residue $k$.
The ring $R$ is regular, excellent and strictly henselian, with
dimension $r+2$. Its quotient by $f$ is the strict henselization of
the cone's vertex ring; its puncture $Y$ is regular. The integral
strict-cone injection C.3e, whose hypotheses and proof apply over
every algebraically closed field, sends $\alpha$ to a non-torsion
class $\gamma\in\mathrm{CH}^r(Y)$. It is the flat pullback to the
punctured line bundle followed by strict localization. Pulling back
the actual proper-curve family above shows
$\gamma\in\mathrm{CH}^r(Y)_{\mathrm{alg}}$ integrally.

Finally, the integral Milnor Bloch formula
\citep[Theorem~7.5, author p.~22]{cyKerz2009Milnor} applies to the
regular noetherian local ring $R$, which contains the infinite field
$k$. It gives
\[
 \mathrm{CH}^r(\operatorname{Spec}R)
       =H^r_{\mathrm{Zar}}(\operatorname{Spec}R,\mathcal K^M_r)=0.
\]
For any abelian sheaf on a local scheme, global sections are its
stalk at the closed point: that point has no proper open
neighbourhood. Global sections are therefore exact, which explains
the displayed higher-cohomology vanishing without an assumption
of finite type for $R$. Chow localization surjects this zero group
onto $\mathrm{CH}^r(X)$. This proves the asserted failure of ordinary
restriction surjectivity.
\end{proof}

The retained exact program checks nonzero Hasse coefficients and unit
Jacobian ideals in every projective chart for specified small-prime
examples. It records when a prime-field parameter search is exhausted;
that is not a failed universal claim, because the construction explicitly
allows a finite extension. The universal existence, rational
non-torsion and actual-family assertions are proved above, not inferred
from a finite calculation.
